\section{Tree Grammar: Domain IR}
\label{sec:tree-grammar}

The Tree Grammar is the fourth grammar in the diagram compiler and the \textbf{second de-risked grammar}: its layout is the Buchheim--J{\"u}nger--Leipert tidy-tree algorithm~\cite{buchheim2002}---deterministic, $O(n)$, and fully solved. Where the Flow Grammar (Section~\ref{sec:flow-grammar}) requires four-phase Sugiyama layered layout with cycle removal and crossing minimization (Section~\ref{sec:layout-engines}), and the Sequence Grammar (Section~\ref{sec:sequence-grammar}) sidesteps layout entirely via arithmetic placement, the Tree Grammar occupies a middle ground: it \emph{does} require a real layout algorithm, but one that is provably deterministic and linear-time for its constrained input class (rooted trees).

% ============================================================================
\subsection{Scope \& Intent}
\label{sec:tree-scope}
% ============================================================================

\subsubsection{What a Tree Diagram IS}

A \textbf{tree diagram} in this grammar is a \emph{node-link rendering of a rooted tree}: a connected, acyclic graph with exactly one root node and where every non-root node has exactly one parent. Trees represent containment, decomposition, and hierarchical relationships:

\begin{itemize}
  \item \textbf{Document structure}: document $\rightarrow$ chapters $\rightarrow$ sections $\rightarrow$ subsections.
  \item \textbf{Organizational hierarchies}: CEO $\rightarrow$ VPs $\rightarrow$ directors $\rightarrow$ teams.
  \item \textbf{Decision trees}: root question $\rightarrow$ branches at each decision point.
  \item \textbf{File systems}: root directory $\rightarrow$ subdirectories $\rightarrow$ files.
  \item \textbf{Taxonomies and classification}: kingdom $\rightarrow$ phylum $\rightarrow$ class $\rightarrow$ order.
\end{itemize}

The visual form is a \emph{tidy tree}: nodes placed at discrete depth levels with sibling groups spread horizontally, connected by parent-to-child edges. This is the canonical ``org chart'' or ``tree view'' layout.

\subsubsection{What a Tree Diagram is NOT}

\begin{description}
  \item[Not a general directed graph.] Flows (Section~\ref{sec:flow-grammar}) support cycles, multi-parent nodes, and arbitrary DAG topologies. Trees require exactly one parent per non-root node---no shared children, no cycles.
  \item[Not a sequence diagram.] Sequence diagrams (Section~\ref{sec:sequence-grammar}) have a rigid participant$\times$message grid structure. Trees have variable branching and depth.
  \item[Not a timeline.] Timelines have a calibrated time axis. Tree depth levels are ordinal (level 0, 1, 2\ldots), not temporal.
  \item[Not a mind map.] Mind maps are typically radial, free-form, and allow cross-links. This grammar produces strictly hierarchical tidy-tree layouts with no cross-edges.
\end{description}

\subsubsection{Modeling Boundary}

The Tree Grammar covers diagrams where:
\begin{enumerate}
  \item Entities are \textbf{nodes} in a containment/decomposition hierarchy.
  \item Relationships are strictly \textbf{parent$\rightarrow$child} (one parent per node, one root).
  \item Layout admits a natural \textbf{level-based reading} (root at top, children below).
\end{enumerate}

Diagrams with shared children (a node having two parents) belong to the Flow grammar (DAG). Diagrams with no clear hierarchy belong to the Graph grammar's force-directed or hub-spoke layouts.

\subsubsection{Why De-Risked}

\textbf{Tidy-tree layout is deterministic and $O(n)$.} The Buchheim--J{\"u}nger--Leipert algorithm~\cite{buchheim2002} (correcting Walker~\cite{walker1990}, itself extending Reingold--Tilford~\cite{reingold1981}) computes compact, non-overlapping node positions in a single linear-time pass:

\begin{itemize}
  \item No NP-hard subproblems (contrast: crossing minimization in Sugiyama is NP-complete~\cite{garey1983}).
  \item No iterative convergence (contrast: force-directed methods).
  \item No randomized initialization.
  \item Fully deterministic given a fixed sibling order.
  \item Proven $O(n)$ time and space.
\end{itemize}

This makes Tree the lowest-risk graph-layout grammar after Sequence (which requires no layout algorithm at all). The ``hard problem'' of Section~\ref{sec:layout-engines} is bypassed by the structural constraint: trees are a strict, easier subclass of general graphs.

% ============================================================================
\subsection{Tree Domain IR}
\label{sec:tree-domain-ir}
% ============================================================================

The Tree Domain IR is the small, grammar-specific representation that compiles \emph{down} to the shared Scene IR (Section~\ref{sec:scene-ir}). It is \textbf{not} a ``god-IR''---it contains only the constructs necessary to describe rooted-tree hierarchies.

\subsubsection{Document Structure}

A Tree IR document is a single YAML/JSON object with the following root structure:

\begin{lstlisting}[language=yaml,caption={Tree IR root document structure}]
version: "1.0"
metadata:
  title: string             # required
  subtitle: string          # optional; default null
  theme: string             # optional; default "default-tree"
tree:
  root: TreeNode            # required; the single root node
\end{lstlisting}

\subsubsection{Root Fields}

\begin{table}[h]
\centering
\small
\begin{tabular}{llccp{4.8cm}}
\toprule
\textbf{Field} & \textbf{Type} & \textbf{Req} & \textbf{Default} & \textbf{Description} \\
\midrule
\texttt{version} & \texttt{string} & Yes & --- & Spec version (semver, e.g., \texttt{"1.0"}) \\
\texttt{metadata} & \texttt{TreeMetadata} & Yes & --- & Document-level metadata \\
\texttt{tree} & \texttt{TreeDefinition} & Yes & --- & The tree definition (contains the root) \\
\bottomrule
\end{tabular}
\caption{Tree IR root fields}
\label{tab:tree-root-fields}
\end{table}

\subsubsection{TreeNode}
\label{sec:tree-node}

A \texttt{TreeNode} is the sole structural construct. Trees are recursive: each node may contain child nodes.

\begin{table}[h]
\centering
\small
\begin{tabular}{llccp{4.0cm}}
\toprule
\textbf{Field} & \textbf{Type} & \textbf{Req} & \textbf{Default} & \textbf{Description} \\
\midrule
\texttt{id} & \texttt{id} & Yes & --- & Document-unique identifier (kebab-case) \\
\texttt{label} & \texttt{string} & Yes & --- & Display text for the node \\
\texttt{children} & \texttt{list<TreeNode>} & No & \texttt{[]} & Ordered child nodes \\
\texttt{kind} & \texttt{string?} & No & \texttt{null} & Semantic kind hint (e.g., \texttt{"chapter"}, \texttt{"person"}, \texttt{"decision"}) \\
\texttt{icon} & \texttt{string?} & No & \texttt{null} & Icon registry key \\
\texttt{collapsed} & \texttt{bool} & No & \texttt{false} & Hint: render subtree as collapsed (elided) \\
\texttt{description} & \texttt{string?} & No & \texttt{null} & Tooltip or secondary text \\
\bottomrule
\end{tabular}
\caption{TreeNode fields}
\label{tab:tree-node-fields}
\end{table}

\paragraph{Canonical representation: children-list.}
The \textbf{canonical} tree representation uses an embedded \texttt{children} list on each node. This is chosen over a flat-list-with-parent-ref representation for three reasons:
\begin{enumerate}
  \item \textbf{Natural authoring}: trees are written top-down; nesting children inside parents mirrors the mental model.
  \item \textbf{Structural guarantee}: a well-formed nested document is automatically a valid tree (no cycles, no orphans, exactly one root). Validation is trivial---only duplicate \texttt{id} checks are needed.
  \item \textbf{Sibling order is explicit}: the list order of \texttt{children} defines left-to-right placement. No separate ordering field is required.
\end{enumerate}

\paragraph{Alternative: flat list with parent refs.}
An alternative flat representation (a list of nodes where each non-root carries a \texttt{parent} ref) is a valid \emph{input convenience}. If supported, the validator \textbf{normalizes} it to the canonical children-list form before layout:
\begin{enumerate}
  \item Verify exactly one node has no \texttt{parent} (the root).
  \item Verify all \texttt{parent} refs resolve to existing node ids.
  \item Verify no cycles (topological sort succeeds).
  \item Group nodes by parent; sibling order = list order among nodes sharing the same parent.
\end{enumerate}

\textbf{Open question for Mark}: Whether to support the flat parent-ref form as an alternative input mode, or require the canonical children-list form exclusively. The spec defines the canonical form; the schema may optionally accept both.

\paragraph{Identity.}
Node \texttt{id} is the primary key across the entire document. It must be globally unique (across all levels of nesting). Two nodes with the same \texttt{id}---regardless of depth---are a validation error.

\paragraph{Sibling order.}
Children are placed left-to-right in the order they appear in the \texttt{children} list. This order is the \textbf{canonical horizontal order} and the sole determinant of sibling arrangement. Reordering children in the IR changes the diagram (this is intentional; order is semantic).

\paragraph{Kind, icon, collapsed: semantic hints only.}
Per the grammar$=$semantics / theme$=$style principle (Section~\ref{sec:grammar-concept}):
\begin{itemize}
  \item \texttt{kind} is a semantic tag (``chapter'', ``person'', ``folder'') that a \texttt{TreeTheme} \emph{may} use to select visual treatment (e.g., different fill colors or border styles per kind). The grammar assigns no visual meaning to kind values---that is a theme concern.
  \item \texttt{icon} references the shared icon registry. Placement and sizing are theme-driven.
  \item \texttt{collapsed} is a rendering hint: when \texttt{true}, the subtree rooted at this node is visually elided (e.g., replaced by an expander indicator). The node itself is still rendered; its children are hidden. This is a \emph{hint}---static backends render it as a visual marker; interactive backends may allow expansion.
\end{itemize}

None of these fields affect layout geometry of \emph{other} nodes. A collapsed subtree occupies zero layout space for its hidden children (the collapsed node itself retains its measured size).

% ============================================================================
\subsection{Layout Contract (Deterministic Tidy-Tree)}
\label{sec:tree-layout-contract}
% ============================================================================

The Tree Grammar's layout engine converts the Domain IR into Scene IR coordinates using the layered tidy-tree algorithm family. The layout is \textbf{fully deterministic}: identical IR in $\rightarrow$ byte-identical Scene out.

\subsubsection{Algorithm: Buchheim--J{\"u}nger--Leipert (2002)}

The layout follows the Reingold--Tilford~\cite{reingold1981} / Walker~\cite{walker1990} / Buchheim--J{\"u}nger--Leipert~\cite{buchheim2002} lineage:

\begin{description}
  \item[Depth assignment.] Each node is assigned to a \textbf{depth level} $d$: the root at $d=0$, its children at $d=1$, and so on. Depth determines the y-coordinate: $y = d \times \text{levelHeight}$.
  
  \item[Horizontal placement.] Siblings are spread horizontally such that:
  \begin{itemize}
    \item No two node bounding boxes overlap (minimum horizontal separation = \texttt{siblingGap}).
    \item Subtrees of adjacent siblings do not overlap (the algorithm walks contours of left and right subtrees, advancing a ``thread'' pointer to detect conflicts).
    \item A parent is centered above its children (the midpoint of the leftmost and rightmost child x-positions).
  \end{itemize}
  
  \item[Complexity.] $O(n)$ time and $O(n)$ space, where $n$ is the number of nodes~\cite{buchheim2002}. The thread-pointer technique avoids Walker's~\cite{walker1990} quadratic worst case.
  
  \item[Determinism.] The algorithm is a pure function of the tree structure and sibling order. No randomness, no iteration count, no convergence criterion. Same input $\rightarrow$ same output, always.
\end{description}

\subsubsection{Orientation}

The primary (default) orientation is \textbf{top-down}: root at the top of the canvas, children below. The depth axis maps to y (increasing downward); the sibling axis maps to x.

\textbf{Left-to-right} orientation (root at left, children to the right) is supported as a theme/layout option---it is the same algorithm with axes transposed. Additional orientations (bottom-up, right-to-left) follow the same transposition pattern.

\textbf{Note}: Orientation is a layout/theme option, not a grammar-level semantic field. The IR describes structure; orientation is presentation. A \texttt{TreeTheme} declares which orientation to use.

\subsubsection{Node Sizing}

Each node's bounding box is computed as:
$$
  w_i = \text{measureText}(\text{node}_i.\text{label}) + 2 \times \text{nodePadX} + \text{iconWidth}_i
$$
$$
  h_i = \text{lineHeight} + 2 \times \text{nodePadY}
$$

All sizing parameters (\texttt{nodePadX}, \texttt{nodePadY}, \texttt{iconWidth}, \texttt{lineHeight}) are \textbf{theme tokens}. The grammar does not prescribe pixel values---only the formula. This ensures that node size is a pure function of label content + theme geometry.

\subsubsection{Spacing Parameters (Theme Tokens)}

\begin{table}[h]
\centering
\small
\begin{tabular}{lp{7cm}}
\toprule
\textbf{Token} & \textbf{Meaning} \\
\midrule
\texttt{levelHeight} & Vertical distance between depth levels (center-to-center) \\
\texttt{siblingGap} & Minimum horizontal gap between adjacent sibling bounding boxes \\
\texttt{subtreeGap} & Minimum horizontal gap between adjacent subtrees (may differ from \texttt{siblingGap}) \\
\texttt{nodePadX}, \texttt{nodePadY} & Internal padding within a node box \\
\texttt{marginTop}, \texttt{marginLeft} & Canvas margins \\
\bottomrule
\end{tabular}
\caption{Tree layout theme tokens (geometry)}
\label{tab:tree-layout-tokens}
\end{table}

\subsubsection{Contrast with Flow Grammar's Sugiyama}

\begin{center}
\small
\begin{tabular}{lp{5cm}p{5cm}}
\toprule
\textbf{Aspect} & \textbf{Flow Grammar (Sugiyama)} & \textbf{Tree Grammar (Tidy-Tree)} \\
\midrule
Input class & General directed graphs (DAGs, cyclic with FAS removal) & Rooted trees only (one parent per node, no cycles) \\
Phases & 4 (cycle removal, layer assignment, crossing min., coordinate assignment) & 1 (recursive bottom-up + top-down adjustment) \\
Crossing minimization & NP-complete; barycenter heuristic, 24 sweeps & Not applicable---trees have no edge crossings \\
Complexity & $O(|V|^2)$ for crossing minimization & $O(n)$ \\
Determinism mechanism & Fixed sweep count + canonical tie-breaking & Inherently deterministic (pure recursion over fixed structure) \\
\bottomrule
\end{tabular}
\end{center}

Trees are a strict subclass of DAGs. The tidy-tree algorithm exploits this constraint to eliminate all hard subproblems that Sugiyama must address.

% ============================================================================
\subsection{Lowering to the Scene IR}
\label{sec:tree-lowering}
% ============================================================================

The tree layout engine emits Scene IR primitives (Section~\ref{sec:scene-ir}). \textbf{No new Scene IR primitives are required}---the existing kernel is sufficient.

\subsubsection{Node Lowering}

Each \texttt{TreeNode} emits a \texttt{Group} containing:

\begin{enumerate}
  \item A \texttt{Rect} primitive for the node box. Shape is theme-driven: default is rounded rectangle (\texttt{corner\_radius: theme.tree.nodeRadius}). The theme may select different shapes per \texttt{kind} (e.g., circles for leaf nodes, diamonds for decision nodes).
  \item A \texttt{Text} primitive for the node label, centered within the box.
  \item Optionally, an \texttt{Image} primitive for the icon (positioned left of label or above, per theme).
\end{enumerate}

\paragraph{Collapsed nodes.}
When \texttt{collapsed: true}, the node is rendered normally but:
\begin{itemize}
  \item Its children are \emph{not} emitted to the Scene IR (zero primitives for the subtree).
  \item An expander indicator is added: a small \texttt{Path} or \texttt{Text} glyph (e.g., ``$+$'' or ``\ldots'') positioned below or beside the node, per theme.
\end{itemize}

\subsubsection{Edge Lowering}

Each parent$\rightarrow$child relationship emits a connector:

\begin{description}
  \item[Orthogonal elbow (default).] A \texttt{Path} with three segments: vertical from parent bottom-center down to the midpoint between levels, horizontal to the child's x-position, then vertical down to the child's top-center. This produces the classic ``org chart'' connector style.
  
  \item[Straight line.] A single \texttt{Line} from parent bottom-center to child top-center. Simpler but less readable for wide trees.
  
  \item[Curved.] A cubic B\'ezier \texttt{Path} from parent to child. Produces smooth, flowing connectors.
\end{description}

The choice among these styles is a \textbf{theme option} (\texttt{theme.tree.edgeStyle: elbow | straight | curved}). The grammar does not prescribe edge rendering---only that a visual connector exists for each parent-child pair.

\paragraph{Edge directionality.}
Edges are rendered \textbf{without arrowheads} by default (the top-down reading direction is implicit). Arrowheads are a theme option for cases where directionality needs emphasis.

\subsubsection{Scene IR Mapping Summary}

\begin{center}
\small
\begin{tabular}{lp{8cm}}
\toprule
\textbf{Tree Construct} & \textbf{Scene IR Primitive(s)} \\
\midrule
Node box & \texttt{Rect} (rounded, theme-styled) \\
Node label & \texttt{Text} (centered in box) \\
Node icon & \texttt{Image} (optional, positioned per theme) \\
Parent$\rightarrow$child edge & \texttt{Path} (elbow/straight/curved, per theme) \\
Collapsed indicator & \texttt{Path} or \texttt{Text} (expander glyph) \\
Canvas & Computed from tree bounding box + margins \\
\bottomrule
\end{tabular}
\end{center}

\paragraph{No new kernel primitives needed.}
The entire Tree Grammar maps to existing Scene IR primitives: \texttt{Rect}, \texttt{Text}, \texttt{Path}, \texttt{Line}, \texttt{Image}, and \texttt{Group}. No kernel changes are required. If future extensions (e.g., cross-links between tree nodes for showing relationships outside the hierarchy) prove necessary, those would require kernel-level review---\textbf{flagged for Barbara/Mark sign-off} rather than assumed.

\paragraph{Styling deferred to TreeTheme.}
Consistent with the \texttt{SequenceTheme} and \texttt{FlowTheme} precedent, all visual styling (colors, font sizes, border widths, edge styles, node shapes per kind, spacing) is controlled by a \texttt{TreeTheme} type. The grammar specifies structure; the theme specifies presentation.

% ============================================================================
\subsection{Edge Cases \& Total Semantics}
\label{sec:tree-edge-cases}
% ============================================================================

Every possible input must have a defined meaning or produce a clear validation error.

\subsubsection{Multiple Roots (Forest)}

\textbf{Rejected.} The Tree Grammar requires exactly one root. A document with multiple roots (e.g., a flat list with two nodes having no parent) is a validation error: \texttt{"Tree must have exactly one root node; found N root candidates."}

\textbf{Rationale}: A forest (multiple disjoint trees) is a fundamentally different layout problem---it requires deciding how to arrange separate trees relative to each other. This is a composition concern. Authors who need a forest should use the Composition layer (Section~\ref{sec:composition}) with multiple Tree panels.

\subsubsection{Cycles}

\textbf{Rejected.} A cycle in the tree structure (node $A$ is a descendant of node $B$ which is a descendant of $A$) is a validation error: \texttt{"Cycle detected involving node 'X'; tree structures must be acyclic."}

In the canonical children-list representation, cycles are structurally impossible (a node cannot appear as its own ancestor in a nested JSON/YAML document). Cycles can only arise in the flat parent-ref alternative input form, where validation must perform a topological sort to detect them.

\subsubsection{Orphan Nodes (Unresolved Parent)}

\textbf{Rejected.} In the flat parent-ref input form, if a node references a parent \texttt{id} that does not exist, validation fails: \texttt{"Node 'X' references non-existent parent 'Y'."}

In the canonical children-list form, orphans are structurally impossible---every node is either the root or a child of another node.

\subsubsection{Single-Node Tree}

\textbf{Valid.} A tree with only a root node (no children) is valid. It produces a single node box centered on the canvas. This is the minimal valid tree.

\subsubsection{Very Wide Trees (Many Siblings)}

\textbf{Valid.} A node with many children (e.g., 50+ siblings) produces a wide diagram. The layout engine does not impose a hard limit on sibling count. The tidy-tree algorithm handles this correctly---it simply places all siblings in a row. A lint warning may be emitted for extreme cases: \texttt{"Node 'X' has N children; consider restructuring for readability."} (Threshold is theme-configurable, e.g., $>20$.)

\subsubsection{Very Deep Trees}

\textbf{Valid.} A tree with many levels (e.g., 20+ depth) produces a tall diagram. No hard limit is imposed. The $O(n)$ algorithm handles deep trees efficiently. As with wide trees, a lint warning may be emitted for extreme depth.

\subsubsection{Duplicate IDs}

\textbf{Validation error.} Two nodes anywhere in the tree with the same \texttt{id} fail validation: \texttt{"Duplicate node id 'X' at paths 'P1' and 'P2'."} (Paths indicate the nesting location for debugging.)

\subsubsection{Empty Children List}

\textbf{Valid.} A node with \texttt{children: []} is a leaf node. This is the common case for terminal nodes and is not degenerate.

% ============================================================================
\subsection{Worked Example: Document Structure Hierarchy}
\label{sec:tree-worked-example}
% ============================================================================

This example models a document hierarchy: a root ``Document'' node with three chapters, each containing a few sections. This is drawn from the corpus's document-structure diagram kind.

\subsubsection{Tree Domain IR}

\begin{lstlisting}[language=yaml,caption={Document hierarchy --- Tree Domain IR}]
version: "1.0"
metadata:
  title: "Document Structure"
  theme: default-tree
tree:
  root:
    id: doc
    label: "Document"
    kind: root
    children:
      - id: ch1
        label: "Chapter 1: Introduction"
        kind: chapter
        children:
          - id: s1-1
            label: "1.1 Background"
            kind: section
          - id: s1-2
            label: "1.2 Motivation"
            kind: section
      - id: ch2
        label: "Chapter 2: Methods"
        kind: chapter
        children:
          - id: s2-1
            label: "2.1 Data Collection"
            kind: section
          - id: s2-2
            label: "2.2 Analysis"
            kind: section
          - id: s2-3
            label: "2.3 Validation"
            kind: section
      - id: ch3
        label: "Chapter 3: Results"
        kind: chapter
        children:
          - id: s3-1
            label: "3.1 Findings"
            kind: section
\end{lstlisting}

\subsubsection{Layout Computation}

The layout engine processes this document as follows:

\begin{enumerate}
  \item \textbf{Depth assignment.} Three levels: \texttt{doc} at $d=0$; \texttt{ch1}, \texttt{ch2}, \texttt{ch3} at $d=1$; all sections at $d=2$.
  
  \item \textbf{Node sizing.} Each node's width is measured from its label text + padding. Suppose (with default theme): \texttt{doc}$=180$px, chapters$\approx220$px, sections$\approx200$px wide; all$=36$px tall.
  
  \item \textbf{Tidy-tree placement (Buchheim et al.).}
  \begin{itemize}
    \item Leaf nodes (sections) are placed with \texttt{siblingGap}$=20$px between them.
    \item Each chapter is centered above its children.
    \item Subtree gaps ensure \texttt{ch1}'s rightmost section and \texttt{ch2}'s leftmost section maintain at least \texttt{subtreeGap}$=40$px separation.
    \item The root \texttt{doc} is centered above the three chapter nodes.
  \end{itemize}
  
  \item \textbf{Vertical positions.} With \texttt{levelHeight}$=100$px: $y_{\text{doc}}=50$, $y_{\text{chapters}}=150$, $y_{\text{sections}}=250$.
\end{enumerate}

\subsubsection{Lowering to Scene IR}

The resulting Scene contains:

\begin{itemize}
  \item \textbf{10 node Groups}: each containing a \texttt{Rect} (rounded, theme-styled) + \texttt{Text} (label). The root node may have a distinct fill color (theme resolves \texttt{kind: root} to a highlight fill).
  
  \item \textbf{9 edge Paths}: one for each parent$\rightarrow$child relationship. With the default \texttt{elbow} edge style, each path has three segments (vertical-horizontal-vertical).
  
  \item \textbf{Canvas}: width computed from the leftmost section's left edge to the rightmost section's right edge, plus margins. Height computed from top margin to depth-2 bottom edge, plus footer margin.
  
  \item \textbf{Meta}: \texttt{scene\_hash} computed from the serialized element list.
\end{itemize}

\paragraph{Primitive count.}
Total Scene primitives: 10 Groups (containing 10 Rects + 10 Texts) + 9 Paths (edges) + canvas + meta = 29 drawing primitives. All are existing kernel types---no extensions required.

\paragraph{Determinism verification.}
Running this IR through the Buchheim layout engine twice produces an identical \texttt{scene\_hash}. Node positions are a pure function of the tree structure, sibling order, and theme geometry tokens. No randomness, no iteration, no heuristic.

% ============================================================================
\subsection{Determinism \& Theme/Semantics Split}
\label{sec:tree-determinism}
% ============================================================================

The Tree Grammar respects the project's sacred determinism contract (Section~\ref{sec:determinism}) and the grammar$=$semantics / theme$=$style principle:

\begin{description}
  \item[Byte-identical output.] The same Tree IR + same engine version + same theme = identical Scene IR (and therefore identical SVG). No exceptions.
  
  \item[No-change defaults.] All optional fields (\texttt{kind}, \texttt{icon}, \texttt{collapsed}, \texttt{description}, \texttt{children}) have explicit defaults. Omitting optional fields produces the same output as specifying the default value explicitly.
  
  \item[New tokens are opt-in.] Future additions to the Tree IR (e.g., new semantic hint fields, annotation support) must have no-change defaults: existing documents that do not use new features must produce byte-identical output after an engine upgrade.
  
  \item[Canonical ordering.] Sibling order in \texttt{children} lists is the sole determinant of horizontal arrangement. Reordering children changes the diagram (intentional---order is semantic).
  
  \item[Deterministic algorithm.] Buchheim--J{\"u}nger--Leipert is a pure function over a tree structure. No iteration count, no sweep count, no RNG. This provides the strongest possible determinism guarantee short of trivial arithmetic (Sequence Grammar).
  
  \item[Grammar is semantic-only.] The Tree Domain IR carries \emph{structure} (parent-child hierarchy, sibling order) and \emph{semantic hints} (\texttt{kind}, \texttt{icon}, \texttt{collapsed}). It does \textbf{not} carry:
  \begin{itemize}
    \item Colors, fills, or strokes (theme tokens).
    \item Font sizes or families (theme tokens).
    \item Node shapes (theme resolves \texttt{kind} to a shape).
    \item Edge styles (theme option: elbow/straight/curved).
    \item Spacing values (theme geometry tokens).
    \item Orientation (theme layout option).
  \end{itemize}
  All visual presentation is controlled by a \texttt{TreeTheme} type, following the \texttt{SequenceTheme} and \texttt{FlowTheme} precedent.
\end{description}

% ============================================================================
\subsection{Open Questions (Deferred to Mark/Barbara)}
\label{sec:tree-open-questions}
% ============================================================================

The following design points are flagged for specialist review before the schema is finalized:

\begin{description}
  \item[Mark (IR schema detail):]
  \begin{itemize}
    \item \textbf{Canonical form}: Should the schema accept \emph{only} the children-list form, or also support a flat parent-ref alternative input? If both, specify the normalization algorithm and which takes precedence on conflict.
    \item \textbf{Forest handling}: Confirm rejection of multiple roots. If future requirements demand forest support, specify it as a separate grammar or as a composition-layer concern.
    \item \textbf{Validation invariants}: Exhaustive list of semantic checks beyond schema conformance (duplicate id detection across nested levels, maximum depth/breadth lint thresholds).
    \item \textbf{Kind enum}: Should \texttt{kind} be a free string or a closed enum? Free string is more extensible (themes can map arbitrary kinds); closed enum enables schema-level validation.
    \item \textbf{Node \texttt{id} format}: Confirm kebab-case requirement. Nested ids may benefit from path-based namespacing (e.g., \texttt{ch1/s1-1}) vs.\ global flat namespace.
  \end{itemize}
  
  \item[Barbara (rendering semantics):]
  \begin{itemize}
    \item \textbf{Edge routing style}: Default elbow geometry---exact midpoint calculation, corner radius for rounded elbows, minimum segment length before defaulting to straight.
    \item \textbf{Collapsed-node handling}: Visual indicator design (``$+$'' glyph, ellipsis, count badge showing hidden children). Interactive vs.\ static rendering semantics.
    \item \textbf{TreeTheme token surface}: Complete list of theme tokens (colors, typography, geometry, edge style) needed for the initial default theme and at least one showcase theme.
    \item \textbf{Kind$\rightarrow$shape mapping}: Default visual mappings for common kind values (e.g., \texttt{person}$\rightarrow$circle, \texttt{folder}$\rightarrow$rounded-rect with icon). How many built-in kind mappings should the default theme provide?
    \item \textbf{Label overflow}: Behavior when label text exceeds available node width---truncate with ellipsis, wrap to multiple lines, or auto-expand node width (current spec assumes auto-expand via \texttt{measureText}).
  \end{itemize}
\end{description}

\paragraph{Kernel change flag.}
No kernel (Scene IR) changes are required for the Tree Grammar as specified. If future extensions (cross-links between tree nodes, animated expand/collapse transitions, or multi-root forest layout) prove necessary, those would require kernel-level review needing Barbara and Mark sign-off before adoption.
